\documentclass[11pt]{article}
\usepackage[T1]{fontenc}
\usepackage{lmodern}
\usepackage[margin=0.85in]{geometry}
\usepackage{amsmath,amssymb,mathtools,booktabs,longtable,array,enumitem,microtype,xcolor,xurl,hyperref,fancyhdr,needspace}
\hypersetup{colorlinks=true,linkcolor=blue!45!black,urlcolor=blue!45!black,
 pdftitle={All-orders perturbative matching in the November Yang--Mills submission},
 pdfauthor={Working draft for the Yang--Mills review group}}
\setlength{\parindent}{0pt}
\setlength{\parskip}{5pt plus 1pt}
\setlength{\headheight}{14pt}
\setlist{leftmargin=1.5em,itemsep=2pt,topsep=3pt}
\pagestyle{fancy}\fancyhf{}
\fancyhead[L]{\small All-orders perturbative matching}
\fancyhead[R]{\small Draft 0.7}
\fancyfoot[C]{\thepage}
\newcommand{\R}{\mathbb R}
\newcommand{\MS}{\overline{\mathrm{MS}}}
\newcommand{\tr}{\operatorname{tr}}
\newcommand{\Tr}{\operatorname{Tr}}
\newcommand{\ad}{\operatorname{ad}}
\newcommand{\code}[1]{{\small\nolinkurl{#1}}}
\begin{document}
\begin{center}
{\LARGE\bfseries All-orders perturbative matching}\par
\vspace{4pt}
The November 4 OpenAI Yang--Mills submission\\
Draft 0.7 \quad\textbullet\quad 16 November 2026
\end{center}

The November matching theorem gives a short-distance expansion of the
Wilson construction's correlations through every finite order, with bounds
on the remainder. Its coefficients are those of a specified continuum
perturbation theory. The proof compares
lattice and continuum flow equations after local counterterms have been
tuned. The comparison concerns the complete coefficient at each order;
it does not require a correspondence between individual graphs. This note
states the result, explains the proof, and examines the conventions through
two ghost-loop comparisons and a gluon-bubble benchmark.

Our human review is incomplete. The submitted Matching theorem has passed
the recorded formal checks, and a separate check verifies its quantifier
structure using the submitted reference. The finite two-point ghost
contribution has also been checked in Lean and Comparator. The four-leg
ghost-box comparison is worked out below, including the source's scale
sectors, dimensional continuation and finite constant. Lean checks its
source histories and subtraction structure; the final finite-value equality
is an analytic derivation, not yet a Lean theorem. The gluon bubble's full finite-value equality has now also been proved in Lean and checked by Comparator, including the finite contribution from the continued vector trace. Source references
and verification records are given
in Appendix~\ref{sec:source}.

The present review is a limited assessment based on source inspection and
selected independent checks. Within the initial two-week review period we
do not claim exhaustive confirmation of the subtraction and composite-operator
conventions. Section~\ref{sec:review-plan} prioritizes further spot checks;
a full identification with an independently specified conventional reference
is reserved for future work. This revision adds plans and a source inventory,
not new Lean or Comparator results.

\section{The matching theorem}
\label{sec:result}

\subsection{The reference, fields, and coupling}

For a gauge group satisfying the submitted simplicity hypothesis, the public
theorem \code{YangMills.perturbativeMatching} produces a Wilson construction
with a \code{MatchedRun}. One lattice family and one perturbative reference
are chosen before the requested volume sequences. For each strictly
increasing requested sequence, the construction selects an infinite-volume
refinement. Each resulting hierarchy has a calibrated field frame and a
measured coupling, common to all perturbative orders, arities, field degrees
and phases. The frame and coupling may depend on the hierarchy, but are
chosen before the perturbative order.

Let $V_{\le d}$ denote the finite sector of gauge-invariant free on-shell
curvature-jet polynomials of engineering degree at most $d$. A $j$th
derivative of curvature has degree $j+2$. Write
\[
 J_r^{(d)}:V_{\le d}\longrightarrow\mathcal H.\mathrm{Field}(d),
 \qquad g_X(r)>0,\qquad g_X(r)\longrightarrow0\quad(r\downarrow0).
\]
The frame $J_r$ includes engineering rescaling and field renormalization.
For fields $v_i\in V_{\le d_i}$, center $o$, and a separated shape
$x=(x_1,\ldots,x_n)$, set
\[
 S_{\sigma,r}(v;x,o)=
 S_\sigma(J_rv_1,\ldots,J_rv_n;o+rx_1,\ldots,o+rx_n).
\]
Here $\sigma$ includes the mixed hierarchy and its character phases, in the
sense specified by the construction.

\subsection{Correlations}

For every finite $N$, derivative order $j$, tuple of degree cutoffs, and
compact set $K$ of separated shapes, there are $A$ and $r_0>0$ such that
\begin{equation}\label{eq:matching}
 \left\|D_x^\ell\left[
 S_{\sigma,r}(v;x,o)-\sum_{k=0}^{N}g_X(r)^k C^X_{n,k}(v;x)
 \right]\right\|
 \le A\prod_i\|v_i\|_{\rm coeff}\,|g_X(r)|^{N+1}
\end{equation}
for $0<r<r_0$, $x\in K$, $0\le\ell\le j$, all allowed fields, all centers
and all the indicated phases. The constants may depend on $N,j,K$ and the
degree cutoffs, and are uniform in the fields, phase, center and point.
The derivatives act on the shape variables $x$ at fixed $r$.

The $C^X_{n,k}$ are defined from continuum graph coefficients followed by a
fixed conversion to the position-calibrated scheme $X$, independently of
the correlations on the left side. Equation~\eqref{eq:matching} asserts an
asymptotic expansion at every fixed finite order. It gives no bound uniform
in $N$ and makes no claim of convergence or Borel summability. The index $k$
counts coupling powers; a one-loop four-gluon contribution, for example,
occurs at order $g^4$.

\subsection{OPE, stress, and order of limits}

For every finite output cutoff $h$, sum the actual OPE channel fields through
that degree and pull the output back by $J_r^{-1}$. The result
$\mathcal C_r^{\le h}$ satisfies the analogous coefficient-norm estimate
\begin{equation}\label{eq:ope}
 \left\|\mathcal C_r^{\le h}(v;x,o)
 -\sum_{k=0}^N g_X(r)^k C_{\mathrm{OPE},k}^{X,\le h}(v;x)\right\|_{\rm coeff}
 \le A\prod_i\|v_i\|_{\rm coeff}|g_X(r)|^{N+1}.
\end{equation}
The full formulation also controls shape derivatives by testing against the
dual unit ball, uniformly in the center. The output degree $h$ and coupling
order $N$ are independent quantifiers. The channels form a basis of the
whole finite output space. The reference OPE is uniquely characterized by
tests against arbitrary exterior insertions. This OPE clause has no phase
index of its own.

The stress tensor is compared with the normalized Ward stress series. Its
remainder, after engineering rescaling, inverse-frame transport and testing
against linear functionals, is $o(g_X^N)$. The stress clause does not include
the derivative estimate~\eqref{eq:matching}. The full theorem also contains
smeared separated tests and bounds on phase differences by powers of $r$,
consistent with agreement to all fixed coupling orders.

The construction first takes its lattice ultraviolet limit at fixed period,
then a selected large-volume limit. The limit $r\downarrow0$ in
\eqref{eq:matching} is taken within that continuum hierarchy. Dimensional
regulator removal and infrared-cutoff removal in the graph definition are
separate, sequential operations. The theorem does not itself identify
contact terms or prove uniqueness of all nonperturbative volume limits.

\section{Reference perturbation theory}
\label{sec:conventions}

The intended conventional reference is pure Euclidean Yang--Mills theory
in Landau gauge, with Faddeev--Popov ghosts and dimensional minimal
subtraction in the scale convention below. The comparison with the submitted
definitions is summarized in Section~\ref{sec:dictionary}.

\subsection{Action and Feynman rules}

Use the submission's invariant color inner product, an orthonormal real
basis, and $[T_b,T_c]=f^{bc}{}_{a}T_a$. This fixes the color normalization,
including $C_A$. Let $g$ be the formal scalar coupling, before conversion
to $g_X(r)$. Set
\begin{align}
 F^a_{\mu\nu}&=F^a_{0,\mu\nu}+gQ^a_{\mu\nu},&
 F^a_{0,\mu\nu}&=\partial_\mu A^a_\nu-\partial_\nu A^a_\mu,\nonumber\\
 Q^a_{\mu\nu}&=f^{bc}{}_{a}A^b_\mu A^c_\nu,&
 (D_\mu c)^a&=\partial_\mu c^a+g f^{bc}{}_{a}A^b_\mu c^c.
\end{align}
For positive $\xi$, the Euclidean action fixing our signs is
\begin{equation}\label{eq:action}
 S_\xi=\int d^4x\left\{\frac14 F^2+
 \frac{(\partial\cdot A)^2}{2\xi}
 +(\partial_\mu\bar c^a)(D_\mu c)^a\right\},\qquad\text{weight }e^{-S_\xi}.
\end{equation}
The reference is defined using the transverse covariance at $\xi=0$, without
taking an interacting $\xi\downarrow0$ limit. Fourier derivatives are
$ip_\mu$; external momenta are incoming and sum to zero.
For $p\ne0$ the free rules are
\begin{align}
 C^{ab}_{\mu\nu}(p)&=\delta^{ab}
 \left(\frac{\delta_{\mu\nu}}{p^2}-\frac{p_\mu p_\nu}{(p^2)^2}\right),
 \label{eq:gauge-line}\\
 \langle c^a(p)\bar c^b(-p)\rangle_0&=\frac{\delta^{ab}}{p^2},&
 \langle\bar c^a(p)c^b(-p)\rangle_0&=-\frac{\delta^{ab}}{p^2}.
 \label{eq:ghost-lines}
\end{align}
The interaction \emph{log density}, which has the opposite sign to the
interaction action, is
\begin{equation}\label{eq:vertices}
 V_3=-\frac g2\int F_0Q,\qquad
 V_4=-\frac{g^2}{4}\int Q^2,\qquad
 V_{\bar cAc}=-g\int(\partial_\mu\bar c^a)f^{bc}{}_{a}A^b_\mu c^c.
\end{equation}
The vertex tensors are functional derivatives of these polynomials. With
external order $(\bar c^a(q),A^b_\nu,c^c)$ the ghost vertex is
$-igq_\nu f^{bc}{}_{a}$. After differentiation, no further vertex factorial
is appended. Expanding $\exp(V_3+V_4+V_{\bar cAc})$ and using ordered Wick
contractions determines multiplicities, symmetry weights and odd-field
signs. Vacuum components cancel in normalized correlations.

For the gauge-fixed path integral and explicit gauge and ghost rules, see
Srednicki, Sections 71--72, especially equations (72.7)--(72.9)
\cite{srednicki}; the original gauge-fixing construction is due to
Faddeev and Popov \cite{fp}. Srednicki uses Minkowski conventions.
Equations~\eqref{eq:action}--\eqref{eq:vertices} specify our Euclidean
signs and normalizations. These references support our comparison; they
are not citations supplied by OpenAI in the manuscripts or November source
we inspected.

\subsection{Regulator, finite parts, and the measured scheme}

Use
\begin{equation}\label{eq:scheme}
 D=4-2\epsilon,\quad z_{\rm source}=-\epsilon,\quad
 \mu_{\rm DR}^{2\epsilon}\frac{d^Dk}{(2\pi)^D},\quad
 \mu_{\rm DR}^2=\frac{e^{\gamma_E}}{4\pi}\bar\mu^2.
\end{equation}
Subtract negative Laurent powers before taking the finite value. Terms
vanishing with $\epsilon$ in a numerator must be retained until subtraction,
because they can multiply poles. The physical reader uses the scale
coordinate $\log\bar\mu+(\gamma_E-\log4\pi)/2$. The graph definition first
removes the dimensional regulator and then sends the retained infrared
length to infinity. The existence of these limits is established by separate
convergence theorems in the submission.

The coefficients $C^X$ are obtained from the $\MS$ coefficients by a field
and coupling conversion fixed by position-space normalization and a signed
triangle measurement. This conversion is common to the whole hierarchy.
Writing the field conversion schematically as $Z(g)$, OPE transport applies
$Z(g)$ to each input and $Z(g)^{-1}$ to the output. The calibration also fixes
the physical frame used in the matching theorem.

\subsection{Comparison with the submitted definitions}
\label{sec:dictionary}

In the source the classical bulk polynomial is a BV/BRST log density of the
form
\[
 \mathrm{ploy}=-\tfrac14F^2+K\cdot Dc-c^*sc,
 \qquad sc=-\tfrac g2[c,c].
\]
The source identity \code{Class_at} identifies the interaction seeds with
$\mathrm{ploy}_g-\mathrm{ploy}_0$. Restricting $K=-\partial\bar c$ and
$c^*=0$ gives the ghost interaction in \eqref{eq:vertices}. This restriction
is not an invertible identification of the unrestricted BV source space
with scalar-antighost fields. The source identifies its two constructions of continued insertion
polynomials; their comparison with a conventional renormalized operator
system still needs review. Section~\ref{sec:completion} examines every row below.

\begingroup\small
\begin{longtable}{@{}>{\raggedright\arraybackslash}p{0.20\textwidth}>{\raggedright\arraybackslash}p{0.43\textwidth}>{\raggedright\arraybackslash}p{0.32\textwidth}@{}}
\toprule
Convention & Correspondence & Checks and unresolved points\\\midrule
\endfirsthead
\toprule Convention & Correspondence & Checks and unresolved points\\\midrule
\endhead
Color and coupling & Submitted metric, $f^{bc}{}_a$, and $G.\mathrm{casimir}$; conversion to $g_X$ by the fixed calibration. & Trace metric and adjoint Casimir checked; conventional SU($N$) conversion stated in Section~\ref{sec:completion}.\\[4pt]
Gauge line & Twice the source half-line is \eqref{eq:gauge-line} times $1-H(e^{2s}p^2)$, $H(u)=(1+u)e^{-u}$. & Equality checked in Lean, including the half-line factor.\\[4pt]
Ghost orientations and cubic vertex & $K=-\partial\bar c$ gives both signs in \eqref{eq:ghost-lines} and $-igq_\nu f^{bc}{}_a$. & Equalities checked in Lean under this source restriction.\\[4pt]
Gauge vertices & Derivatives of $-gF_0Q/2-g^2Q^2/4$ in the submitted color metric. & Density identities checked in Lean; tensors derived and tested by exact algebra. Generic Lean tensor comparison remains open.\\[4pt]
Subtraction & Pure-pole subtraction in \eqref{eq:scheme}, with child counterterms included in each summed scale sector. & Forest identity, ghost finite value and gluon finite value checked; box compared analytically. Conventional all-order forest/continuation comparison remains open.\\[4pt]
Local fields and OPE & Free on-shell curvature-jet labels, renormalized insertions, common $J_r$, inverse map on the output. & Quantifiers and internal continued-insertion bridge checked. External operator/BRST interpretation still needs review.\\[4pt]
Lattice normalization & Wilson family and local tuning with the same physical calibration as the reference. & Public/internal Wilson law identified; leading normalization derived and tested. General tuning and estimates need further audit.\\
\bottomrule
\end{longtable}
\endgroup
\section{Proof strategy}
\label{sec:proof}

The proof first identifies continuum graph sums with a solution of the
renormalization flow equation. It then compares the lattice coefficient
flow with that solution, and uses finite-order remainder estimates to pass
from perturbative coefficients to the constructed continuum correlations.

\subsection{Continuum flow and uniqueness}

Integrating a momentum shell changes an effective interaction by a term
closing two legs of one kernel and a term joining two kernels by a shell
line. At fixed coupling order the schematic equation is
\begin{equation}\label{eq:flow}
 \partial_t F_n
 =\mathcal L_{\dot C_t}(F_n)
   +\sum_{n_1+n_2=n}\mathcal Q_{\dot C_t}(B_{n_1},F_{n_2}).
\end{equation}
Here $B$ denotes the bulk coefficient family and $F$ either that family or
a family with one distinguished insertion. In the joining term the bulk
order is positive; the insertion coefficient may have order zero. The
bulk joining term carries a factor $1/2$, whereas the term with a
distinguished insertion carries factor $1$. Graded contractions supply the
odd-field signs. The implementation in \code{CoverFlow} uses
\code{cycle}, \code{using_tree} and their sum \code{recu}.

Closing a pair of legs uses a kernel with two more external legs. The
induction therefore proceeds through lower coupling orders and, at a fixed
order, from higher to lower external rank. Finite-rank support and local
boundary data complete the recursion. The flow identity also holds for
momentum derivatives, providing the smoothness and uniform bounds used later.

Continuum graphs are organized as ordered Wick sectors, with internal
assignments and contractions summed. Differentiating with respect to the
shell scale exposes the last contraction or joining operation. The identities
\code{ambient_bob} and
\code{mark_bob} establish the flow for the total bulk and marked graph
coefficients. Pole subtraction and classical input fix the local boundary
data. Uniqueness then identifies these sums with the continuum MS solution:
\code{income_bulk} and \code{income_mark}. Checking the contraction rules,
classical input and boundary conditions is thus the central task in
identifying the reference theory at all orders.

Flow-equation renormalization and composite-operator estimates are developed
in \cite{egk,fhh,fh}, with different conventions. For example, \cite{egk} treats SU(2)
with a cutoff regulator and positive gauge parameter, whereas the present
reference uses Landau gauge and dimensional subtraction.

\subsection{Comparison of lattice and continuum flows}
\label{sec:lattice-flow}

The lattice flow uses the lattice momentum symbol and integrates over the
first Brillouin zone. The continuum flow uses ordinary momentum and
integrates over all Euclidean loop momenta. They have the same contraction
and joining structure, with different shell covariances and loop measures.

Write $F^{\rm lat}_n$ for the rescaled lattice coefficient and
$F^{\rm cont}_n$ for the continuum coefficient. These denote full kernels
at order $n$, including local counterterms. To control
$\Delta_n=F^{\rm lat}_n-F^{\rm cont}_n$, estimates on the two propagators
are combined with the induction hypotheses to bound the difference of the
flow right-hand sides. The same comparison is made for the Ward/master
identity. These estimates are \code{first_rhs} and \code{second_rhs} in
\code{High/Sphere}.

The flow bounds give limits of the low momentum Taylor coefficients of
$\Delta_n$. These local terms are the remaining obstruction to matching.
The Ward/master identity implies that the local discrepancy is BRST closed.
The local classification then expresses it in terms of permitted action
counterterms and, for insertions, field mixing and BRST-exact adjustments.
Their coefficients are tuned to cancel the discrepancy while preserving
all earlier orders. The leading cubic normalization is fixed separately.

Once these Taylor terms vanish, \code{terminal} bounds the entire kernel
difference and its admitted momentum derivatives. The theorems
\code{High.action} and \code{High.mark} choose a compatible sequence of
parameters for all orders.

The formal conclusion, \code{Fixed}, states that every admitted order and
external rank satisfies the weighted estimate \code{Trust}. Its norm
includes engineering scaling, and the difference has a gain
$\exp(-t/2)$, where $t$ is the log heat-length coordinate, distinct from the
physical radius $r$ in the matching theorem. The comparison is applied to the lattice solution by
\code{portH} and \code{leader_M}, using the continuum MS solution as the
reference.

After tuning, the rescaled lattice coefficients therefore approach the
continuum solution of the flow equation. The comparison of the two flows
supplies the control needed to identify this limit. We have traced these steps in
the source; their analytic estimates still require human review.

\subsection{Bounds for physical correlations}

The matched coefficients are transported through the common field and
coupling calibration and the continuum construction. Compatibility of the
local prescriptions, uniqueness of the bulk data, and uniqueness of limits
identify the coefficients obtained along these routes. This later
common-limit argument uses the preceding lattice--continuum flow comparison.

Finite-order expansion estimates for the lattice observables and their
continuum limits then give the remainder in~\eqref{eq:matching}. Momentum
and position derivative estimates and the measured small-coupling law give
the stated uniformities. Insertion estimates and the exterior-field
characterization of the OPE give~\eqref{eq:ope}. The construction, field
frame and measured coupling are the same throughout.

\section{Matching at fixed order}
\label{sec:graphs}

The object in \code{Fixed} is the full coefficient kernel at a coupling
order and external rank. Its arguments contain no graph label. At the
physical level the theorem compares
\[
 \lim_{\text{continuum scaling}}
 \bigl[\hbox{lattice coefficient at order }n\bigr]
 =\hbox{continuum coefficient at order }n,
\]
in the controlled, tuned sense of Section~\ref{sec:lattice-flow}, and then
uses these coefficients in \eqref{eq:matching}. It does not assert an
equality indexed by a lattice graph and a corresponding continuum graph.

The proof uses Ward/BRST identities for sums of contributions: cancellations
between terms can be needed before the local discrepancy is BRST closed.
Counterterms, field mixing and finite scheme changes also redistribute
contributions. Moreover, the lattice action has higher vertices absent
from the classical continuum action; these can disappear in the limit or
contribute to the local terms removed by tuning.

The submitted graph organization reflects this. A \code{Diagram} describes
ordered sectors with sums over internal assignments. Subtraction acts on
integrated sector sums, including child counterterms. A correspondence
between these objects and individual conventional Feynman graphs would
require a further specification of grouping and subtraction.

Graph-by-graph matching may hold under additional choices, but it is not
asserted by the present theorem. The examples below test individual
contributions to the continuum reference, including signs, multiplicities
and finite subtraction constants. Their comparison with the submitted
continuum evaluator does not establish a separate lattice limit for each
graph. They also concern elementary gauge and ghost kernels, whereas the
public matching theorem concerns gauge-invariant composite insertions.

\section{One-loop ghost contribution to the two-point kernel}
\label{sec:bubble}

Strip the overall $g^2$, set $p^2>0$, and use the color contraction
$C^{ab}=C_A\delta^{ab}$ in the submitted metric. The conventional integral
in our log-density conventions is
\begin{equation}\label{eq:bubble-integral}
 B^{ab}_{\mu\nu}(p;\epsilon)
 =-C^{ab}\mu_{\rm DR}^{2\epsilon}
 \int\frac{d^{4-2\epsilon}k}{(2\pi)^{4-2\epsilon}}
 \frac{k_\mu(k+p)_\nu}{k^2(k+p)^2}.
\end{equation}
The sign includes the ordered vertices and the closed odd-field loop.
Introduce one Feynman parameter, shift the loop momentum, and use rotational
tensor reduction in $D$ dimensions. With $L=\log(\bar\mu^2/p^2)$, this gives
\begin{align}
 B^{ab}_{\mu\nu}
 &=\frac{C^{ab}}{16\pi^2}\frac{Q(\epsilon,L)}{\epsilon}
 \left[\frac{\delta_{\mu\nu}p^2}{2(1-\epsilon)}+p_\mu p_\nu\right],
 \label{eq:bubble-gamma}\\
 Q(\epsilon,L)
 &=e^{(L+\gamma_E)\epsilon}
 \frac{\Gamma(1+\epsilon)\Gamma(2-\epsilon)^2}{\Gamma(4-2\epsilon)}.
\end{align}
In particular $Q(0,L)=1/6$ and
$\partial_\epsilon Q(0,L)=(L+5/3)/6$. Subtracting the pole gives
\begin{equation}\label{eq:bubble-answer}
 \boxed{B^{ab}_{\mu\nu,R}(p)=\frac{C^{ab}}{16\pi^2}
 \left[\frac{\delta_{\mu\nu}p^2}{12}\left(L+\frac83\right)
 +\frac{p_\mu p_\nu}{6}\left(L+\frac53\right)\right].}
\end{equation}
The extra finite $1$ in $8/3=5/3+1$ comes from $1/(1-\epsilon)$ multiplying
the pole. It is an example of the finite terms lost if the numerator is
evaluated at $D=4$ before subtraction. This ghost contribution is not
transverse on its own.

The theorem \code{GhostSolution.acceptance} proves
\eqref{eq:bubble-answer} for the actual
\code{(TwoPoint.bubble G true true).instructions.renorm}, read by
\code{Outlook.four} and taken to its infrared limit. It covers every nonzero
momentum, all color and tensor components, and every positive scale under
the stated group hypothesis. The 13 November record contains a fresh
15-module review build, the standard-axiom audit and a successful frozen-target
Comparator run. The formula was specified independently of the new proof.

The proof follows the submitted contraction, regulator derivative, radial
integral, scale flow and endpoint limits, then restores the general tensor
and scale. The conventional integral-to-Gamma reduction
\eqref{eq:bubble-integral}--\eqref{eq:bubble-gamma} is an analytic calculation outside Lean;
the Lean result identifies the submitted graph with the explicit finite
answer. This checks the color normalization, ghost signs, tensor structure,
scale convention and finite subtraction constants for this contribution.
The formal finite-value result in this section concerns the ghost bubble. Section~\ref{sec:gluon-bubble} gives the corresponding gluon calculation and its completed source comparison.

\section{One-loop ghost box with four external gauge legs}
\label{sec:box}

The second calculation compares the finite ghost-box contribution at
specified nonexceptional momenta with the submitted evaluator. The comparison
agrees, including the multiplicity, sign, scale dependence and finite
constant. We give the conventional calculation first, then the source
comparison and the precise scope of its formal checks.

\subsection{Kinematics}

Choose a nonzero Euclidean momentum $p$, a unit polarization $e$ with
$e\cdot p=0$, and incoming momenta
\begin{equation}\label{eq:box-momenta}
 (p_1,p_2,p_3,p_4)=(p,p,p,-3p).
\end{equation}
Every nonempty proper subset has nonzero total momentum. This is therefore
a nonexceptional Euclidean configuration despite being collinear. Take all
four color directions equal to $h$ and write $F_h=\ad_h$, a real
antisymmetric matrix for the fixed invariant metric. The color factor is
$\tr_{\ad}(F_h^4)$.

\subsection{Signs and multiplicity from the ghost determinant}

Integrating the ghosts gives the log-density contribution
\begin{equation}\label{eq:det}
 W_{\rm gh}[A]=\log\det M[A]-\log\det M[0],\qquad M=-\partial_\mu D_\mu.
\end{equation}
Writing $M=M_0+gM_1$ gives
\[
 W_{\rm gh}=\sum_{m\ge1}\frac{(-1)^{m+1}g^m}{m}
 \Tr[(M_0^{-1}M_1)^m].
\]
For this transverse color/polarization projection, each $M_1$ supplies
$-i(e\cdot k)F_h$; shifts of $k$ by multiples of $p$ leave the numerator
unchanged. Four functional derivatives of the $m=4$ term leave
$4!/4=6$ cyclic orderings. They all give the same scalar integral at
\eqref{eq:box-momenta}. Both loop orientations are already included among
these six; there is no additional factor two. Stripping $g^4$ and the
overall momentum delta, the bare projected contribution is
\begin{equation}\label{eq:box-prefactor}
 B^{\rm gh}_{eeee;hhhh}=-6\tr_{\ad}(F_h^4)I(p;\epsilon),
 \quad
 I=\mu_{\rm DR}^{2\epsilon}\int\frac{d^Dk}{(2\pi)^D}
 \frac{(e\cdot k)^4}{k^2(k+p)^2(k+2p)^2(k+3p)^2}.
\end{equation}
This is a log-density sign. The induced effective \emph{action} has the
opposite sign. The integral has no infrared divergence at these momenta;
its one-loop ultraviolet divergence is logarithmic, with no proper loop
subdivergence.

\subsection{Evaluation}

Let $x_0,\ldots,x_3\ge0$, $\sum x_j=1$, and integrate over the flat simplex
$\Delta_3$, whose volume is $1/6$. Define
\[
 V(x)=\sum_{j=0}^3j^2x_j-\left(\sum_{j=0}^3jx_j\right)^2.
\]
Feynman parametrization shifts $k$ by $p\sum jx_j$, leaving $e\cdot k$
unchanged. The rotational average must be done in $D$ dimensions:
\[
 (e\cdot\ell)^4\ \rightsquigarrow\
 \frac{3(\ell^2)^2}{D(D+2)}.
\]
The radial integral and the parameter factorial then give
\begin{equation}\label{eq:box-param}
 I(p;\epsilon)=\frac{3\mu_{\rm DR}^{2\epsilon}}{4(4\pi)^{2-\epsilon}}
 \Gamma(\epsilon)\int_{\Delta_3}(p^2V(x))^{-\epsilon}\,dx.
\end{equation}
Using \eqref{eq:scheme},
\begin{equation}\label{eq:box-expanded}
 I(p;\epsilon)=\frac1{128\pi^2}
 \left[\frac1\epsilon+\log\frac{\bar\mu^2}{p^2}+c_\Box\right]+O(\epsilon),
 \qquad c_\Box=-6\int_{\Delta_3}\log V(x)\,dx.
\end{equation}
At these momenta the parameter integral can be evaluated in elementary
logarithms.

Put $D_j=(k+jp)^2$ and use $D_0-2D_1+D_2=2p^2$ to reduce the four
denominators to three, then reduce each triple in the same way. Grouping
translated pairs by their separation yields
\begin{equation}\label{eq:box-bubbles}
 I(p;\epsilon)=\frac1{(p^2)^2}
 \left[\frac5{12}J(p;\epsilon)-\frac23J(2p;\epsilon)
 +\frac14J(3p;\epsilon)\right],
\end{equation}
where $J(q;\epsilon)$ has numerator $(e\cdot k)^4$ and denominators
$k^2(k+q)^2$, with the same measure as $I$. Shifts along $p$ preserve the
numerator. This reduction uses the same translation-invariant dimensional
continuation for every term; its power-divergent bubbles are not separate
convergent four-dimensional integrals. Their continued values are
\begin{equation}\label{eq:J}
 J(q;\epsilon)=\frac{3\mu_{\rm DR}^{2\epsilon}}{4(4\pi)^{2-\epsilon}}
 \Gamma(\epsilon-2)(q^2)^{2-\epsilon}
 \frac{\Gamma(3-\epsilon)^2}{\Gamma(6-2\epsilon)}.
\end{equation}
The weighted $q^4$ factors in \eqref{eq:box-bubbles} sum to $10$.
Expanding \eqref{eq:J} then gives
\begin{equation}\label{eq:cbox}
 \boxed{c_\Box=\frac{46}{15}+\frac{32}{15}\log2-\frac{81}{20}\log3
 =0.09600088275537\ldots.}
\end{equation}
Pure-pole subtraction and \eqref{eq:box-prefactor} give the finite result
\begin{equation}\label{eq:box-answer}
 \boxed{B^{\rm gh}_{eeee;hhhh,R}
 =-\frac{3\tr_{\ad}(F_h^4)}{64\pi^2}
 \left[\log\frac{\bar\mu^2}{p^2}
 +\frac{46}{15}+\frac{32}{15}\log2-\frac{81}{20}\log3\right].}
\end{equation}
The scalar integral's derivative with respect to $\log\bar\mu^2$ is
$1/(128\pi^2)$; the finite constant additionally tests the subtraction
convention.

\subsection{Identifying and summing the submitted histories}

Let $v$ denote \code{vertex (interaction G 1)}, let $(AB)$ denote
\code{bridge (cursor true) A B}, and close each tree below with
\code{loop (cursor true)}:
\[
 v(v(vv)),\qquad v((vv)v),\qquad (vv)(vv),\qquad
 (v(vv))v,\qquad ((vv)v)v.
\]
These are five ordered tree histories, each with four cubic vertices and
three ghost bridges, followed by one ghost closure. They occur in the
submitted order-four list \code{bulks (bicolor G) 4}. Their membership and
four-leg rank are proved in \code{GhostBoxSource.lean}.
The submitted \code{germ_bulk} and \code{reached_bulk} in
\code{Reaching.lean} identify the full history sum with the MS bulk kernels. We isolate the ghost box
within that representation, before conversion to scheme $X$.

To see why these histories exhaust the chosen ghost-box sector, require
four external gauge legs and only ghost internal lines. The cubic $c^*cc$
vertex cannot occur: every allowed line pairs a $c$ with a $K$, whereas
this vertex adds two unmatched ghosts. A pure gauge cubic cannot join
another vertex through a ghost line. Thus each vertex is $KAc$, with one
external gauge leg and two internal ghost ends. A connected graph is then
a four-vertex cycle. Its latest edge closes the cycle; removing that edge
leaves a tree. No proper subgraph contains a loop. This exhaustiveness
argument is part of the written derivation, not a theorem in the new Lean
file.

The source's \code{bridgeG} supplies twice the half-line and a Koszul
shuffle; the bulk \code{factor true} is $1/2$. A closure sums both ordered
ends against the half-line. The scalar-antighost dictionary in
Section~\ref{sec:dictionary} applies to every internal line, with no external
antifield left over. After factoring out the momentum and color tensors,
transcribing these rules into exact rational arithmetic gives 384 nonzero
contractions per tree, each of weight $-1/16$.
The balanced tree permits either order of its two child bridges. Resolving
that choice gives $2304$ contractions on strictly ordered scale domains.
For each of the six directed cyclic orderings and each of the $4!=24$
edge-scale orders, exactly sixteen terms occur, with total weight $-1$.
Dropping the balanced tree fails this check.

There is also a direct counting explanation. At distinct edge scales, remove
the latest edge of the cycle, then recursively remove the latest edge of
each remaining tree component. This fixes the contraction history. Each of
its three bridges has two child orders and its final closure has two end
orders, giving $2^3\cdot2=16$ terms. The four half-lines supply $2^{-4}$;
the Koszul contraction around the closed ghost cycle supplies the minus
sign. The ordered domains partition the four independent edge-scale
integrals, apart from boundaries of measure zero. Thus the sum reconstructs
the six terms in the ghost determinant, with no additional symmetry factor.

For infrared length $\rho=e^s$, the scalar full line is
\[
 C_\rho(q)=\frac{1-H(\rho^2q^2)}{q^2},\qquad
 H(u)=(1+u)e^{-u},\qquad
 \partial_\rho C_\rho(q)=2\rho^3q^2e^{-\rho^2q^2}.
\]
These are the source identities \code{halfP_components} and
\code{shell_deriv}, after the ghost dictionary. To justify rearrangements,
first retain a positive ultraviolet length $a<\rho$, so each full line is
$C_\rho-C_a$. All edge scales then lie in the compact interval $[a,\rho]$.
After summing sectors and taking the dimensionally continued ultraviolet
limit, the scalar integral is
\begin{equation}\label{eq:box-cutoff}
 I_\rho(p;\epsilon)=\mu_{\rm DR}^{2\epsilon}
 \int\frac{d^Dk}{(2\pi)^D}\,(e\cdot k)^4
 \prod_{j=0}^3\frac{1-H(\rho^2(k+jp)^2)}{(k+jp)^2}.
\end{equation}
The projected ghost contribution is $-6\tr(F_h^4)I_\rho$.

\subsection{Dimensional continuation, subtraction and infrared limit}

Agreement at integer dimensions would not determine the finite part.
Here the source continuation can be checked directly. At a fixed set of
positive edge scales, its split-radial integrand is a polynomial in $r^2$
times $e^{-Ar^2}$, where $A>0$ is the sum of the squared edge lengths.
The ghost line has no vector-polarization trace term \code{yeast}.
For $j=0,\ldots,4$, the source's radial derivative gives
\begin{equation}\label{eq:box-radial}
 -\frac1{\Gamma(1+z)}\int_0^\infty r^{2z}
 \frac{d}{dr}\bigl(r^{2j}e^{-Ar^2}\bigr)\,dr
 =(z)_j A^{-z-j},\qquad \Re z>-1,
\end{equation}
where $(z)_0=1$ and $(z)_j=z(z+1)\cdots(z+j-1)$ for $j>0$.
For $j=0$ this is the elementary Gamma integral; for $j>0$, differentiating
the two terms and using the Gamma recurrence gives the displayed result.
Together with \code{tenorK}, including its $(4\pi)^{-z}$ factor, these are
exactly the extra-$2z$-dimensional Gaussian moments. This identifies the
continuation on an open set, including its first regulator derivative,
with $D=4+2z=4-2\epsilon$. Compact edge-scale ranges justify the intermediate
integrations. The ultraviolet limit then converges for small
$\Re\epsilon>0$ and determines the meromorphic continuation near zero.

All proper trees have zero counterterm. The new Lean theorem
\code{tree_pole} proves this from the submitted \code{ct_blind} and
\code{scaleZero_invariant}; it is not an assumption about the graph.
For each of the five final closures, \code{finite_reading} identifies the
actual \code{Outlook.four} value with \code{oneLoopFinite}. That formula
retains \code{ruby (highTensor ... ) 1}, the derivative of the local tensor
in the dimensional regulator. The only subtraction is therefore the overall
pure pole of \eqref{eq:box-cutoff}. No extra finite counterterm is inserted.

The removal of the infrared cutoff does not introduce a finite shift.
The difference $I-I_\rho$ has no ultraviolet pole: every term contains at
least one exponentially decreasing $H(\rho^2(k+jp)^2)$. It is also infrared
integrable near each of the four distinct soft points. Writing $k=-jp+q$,
the transverse numerator is $(e\cdot q)^4$, so the integrand is bounded
by a constant times $|q|^2$ there. In four dimensions its integral against
the cutoff is $O((\rho|p|)^{-6})$. Away from these points the cutoff gives
exponential decay. The difference is regular near $\epsilon=0$, and hence
\[
 \operatorname{FP}_{\epsilon=0}I_\rho
 =\operatorname{FP}_{\epsilon=0}I-(I-I_\rho)\big|_{D=4}
 \longrightarrow \operatorname{FP}_{\epsilon=0}I.
\]
This proves the stated order of limits: subtract the pole at fixed $\rho$,
then take $\rho\to\infty$. Combining this with the sector sum gives
\eqref{eq:box-answer} for the submitted ghost-box contribution in $\MS$.
The sector reconstruction and integral comparison in this paragraph are
analytic arguments; the final equality has not been formalized in Lean.

\subsection{Independent numerical checks and formal scope}

The accompanying numerical script evaluates the simplex expression for
$c_\Box$ independently of the bubble reduction. Tensor-product Gauss--Legendre
quadrature with 64, 96 and 144 nodes per coordinate differs from
\eqref{eq:cbox} by approximately $4.6\cdot10^{-10}$,
$4.1\cdot10^{-11}$ and $3.6\cdot10^{-12}$, respectively. Exact rational
arithmetic checks the bubble weights and the underlying polynomial
denominator identity; sampled momenta provide an additional numerical
check. The quadrature errors have not been bounded rigorously.

A second calculation integrates the source-cutoff expression directly in
four-dimensional momentum space, without the simplex or bubble reduction.
Set $p^2=\bar\mu^2=1$ and subtract
$(e\cdot k)^4/[(k+3p/2)^2+1]^4$ inside the integral. This massive reference
has the same ultraviolet pole and zero $\MS$ finite part at this scale:
its dimensionally continued integral is
\[
 \frac{\mu_{\rm DR}^{2\epsilon}}{8(4\pi)^{2-\epsilon}}
 \Gamma(\epsilon)
 =\frac1{128\pi^2\epsilon}+O(\epsilon)
 \qquad (\bar\mu=1).
\]
Thus the remaining integral is convergent and equals $I_\rho^{\MS}$.
Integrating the transverse angles analytically leaves two numerical
variables. The resulting values, normalized by $128\pi^2$, are
\begin{center}
\begin{tabular}{@{}rr@{}}
\toprule $\rho$ & $128\pi^2 I_\rho^{\MS}$\\\midrule
4 & 0.095199035282\\
8 & 0.095985284609\\
16 & 0.096000623441\\
32 & 0.096000878608\\
$\infty$ & 0.096000882725\\
\bottomrule
\end{tabular}
\end{center}
The last row evaluates the cutoff-free convergent difference directly. It
differs from the exact constant by $3.1\cdot10^{-11}$; successive quadrature
sizes approach the same value. These numerical errors are not certified
bounds. The calculation checks a finite subtraction constant, rather than
only the pole or scale derivative. For example, prematurely replacing
$3/[D(D+2)]$ by its four-dimensional value in the dimensionally continued
radial integral would shift $c_\Box$ by $-5/6$ and fail this comparison.

The new Lean file passed a fresh build and a standard-axiom audit for six
structural theorems. It also states the exact proposition
\code{GhostBox.SubmittedFinitePartTest}, using the actual five histories
and the independently specified answer, for a future formal proof.
That proposition is not proved, and no new Comparator result is claimed.
The exact sector enumeration is a Python transcription of the source
rules, not execution of the Lean evaluator. Thus the finite comparison is
supported by a written derivation, exact combinatorics and independent
quadrature; only the source identification and subtraction reduction are
currently formalized.

\subsection{Plan to complete the formal ghost-box check}
\label{sec:box-plan}

This is a bounded continuation of the existing benchmark, with lower priority
than the subtraction and operator-dictionary checks in
Section~\ref{sec:review-plan}. Keep \code{GhostBox.SubmittedFinitePartTest}
fixed: arbitrary color direction, positive renormalization scale, nonzero
momentum, and a transverse unit coordinate polarization, at momenta
$(p,p,p,-3p)$. The independently specified right-hand side must not be
redefined using the submitted evaluator.

\begin{enumerate}
\item \textbf{Reconstruct the integral from the actual histories.}
Use \code{tree_result}, \code{tree_pole}, \code{finite_reading} and
\code{actual_pole}. Prove the contraction signs, multiplicities and partition
of the ordered scale domains, including the balanced tree. Lift the exact
sector enumeration into a theorem about the Lean operations, rather than
treating the Python transcription as execution of those operations.
Separately establish exhaustion of the chosen four-cubic-vertex ghost sector:
\code{family_member} proves membership, not completeness.
\item \textbf{Evaluate the dimensionally continued scalar integral.}
Formalize either the simplex calculation or the reduction to the three
bubble integrals at $p,2p,3p$. Keep a common dimensional continuation for
the individual terms, which need not converge separately in four dimensions.
Retain the regulator dependence of the angular factor and Gamma functions
through the finite coefficient; recover the exact rational and logarithmic
constant in \eqref{eq:cbox}.
\item \textbf{Connect subtraction and the infrared limit.}
Identify the actual Laurent pole and finite part in the source convention,
justify the sum/integral rearrangements and cutoff removal, and transport
the result from normalized kinematics to the full quantified target.
Reuse the completed bubble infrastructure where its hypotheses apply.
Numerical convergence is supporting evidence, not a replacement for these
limit estimates.
\item \textbf{Close and audit the target.}
Prove \code{SubmittedFinitePartTest}, freshly build its dependencies, audit
the axioms, and run Comparator against an independently frozen challenge.
Record source-manifest stability and the exact domain of the theorem.
\end{enumerate}

Completion means equality for the submitted finite evaluator, with the
sector identification also stated precisely. It does not establish the
whole four-gluon amplitude. The first milestone is the source-to-integral
identity; its completion will give a firmer estimate of the remaining
formalization effort. No completion date or additional passed check is
claimed here.

\subsection{Scope of the example}

A full one-loop four-gluon check would also include gauge boxes, triangles
with a quartic vertex, bubbles with two quartic vertices, and local
counterterms. In the submitted connected/amputated convention, reducible
assemblies from lower kernels must be included as well. Extracting a 1PI
coefficient would require a separate reconstruction. The comparison must
also specify whether it is made in MS or after conversion to scheme $X$.

For that fuller check, the projection used here should be supplemented by
noncommuting color directions and an independent polarization projection.
Classical commutator vertices can vanish when all colors equal $h$, so this
projection alone does not determine the quartic gauge tensor.

\section{Gluon bubble and the extra vector trace}
\label{sec:gluon-bubble}

The two preceding examples have ghosts in the loop. Here the internal
lines are gluons, so a sum over vector components enters the numerator.
We calculate the conventional Landau-gauge bubble through its finite term
and prove the same finite value for the submitted evaluator. The proof
includes the actual extra-axis trace, radial continuation, finite subtraction
and infrared limit. The earlier trace diagnostics are retained below.

\subsection{Reference integral and conventions}

Use the same Euclidean log-density sign, color metric, Fourier convention
and $\MS$ measure as above. Restore $g^2$ in this section and write
$C^{ab}=C_A\delta^{ab}$ and $L=\log(\bar\mu^2/p^2)$, with $p^2>0$.
Define
\[
 P^T_{\mu\nu}(p)=\delta_{\mu\nu}-\frac{p_\mu p_\nu}{p^2},
 \qquad P^L_{\mu\nu}(p)=\frac{p_\mu p_\nu}{p^2}.
\]
For the two vertices, route the internal momenta as $k$ and $r=k+p$ and
use the real cubic tensor
\[
 V_{\mu\alpha\beta}(p,k,-r)
 =\delta_{\mu\alpha}(p-k)_\beta
  +\delta_{\alpha\beta}(k+r)_\mu
  -\delta_{\beta\mu}(r+p)_\alpha.
\]
The two factors of $i$ and reversal of the momenta at the second vertex
give the positive contraction
\begin{equation}\label{eq:gluon-integral}
 G^{ab}_{\mu\nu}
 =\frac{g^2C^{ab}}2\mu_{\rm DR}^{2\epsilon}
 \int\frac{d^Dk}{(2\pi)^D}
 \frac{V_{\mu\alpha\beta}P^T_{\alpha\alpha'}(k)
 P^T_{\beta\beta'}(r)V_{\nu\alpha'\beta'}}{k^2r^2}.
\end{equation}
The factor $1/2$ is the bubble symmetry factor. This is a log-density
kernel; the induced effective-action kernel has the opposite sign.

To isolate the vector trace, temporarily denote its dimension by $D_s$,
while loop integration and momentum tensor reduction remain in $D$
dimensions. The intended comparison sets $D_s=D=4-2\epsilon$.
Replacing only $D_s$ by $4$ below is a diagnostic modification, not a
specification of a complete alternative regularization scheme.

\subsection{Contraction and reduction}

Set $p^2=1$ temporarily, let $A=k^2$, $B=(k+p)^2$, and take a unit
external polarization $e$. For $e\cdot p=0$, put $t=e\cdot k$. The
contracted numerator in \eqref{eq:gluon-integral}, before the factor $1/2$,
is
\begin{align}
 N_T={}&4+A+B-\frac{(B-1)^2}{A}-\frac{(A-1)^2}{B}\notag\\
 &+t^2\left[4D_s-6+
 \frac{A^2+B^2-2A-2B+1}{AB}\right].\label{eq:gluon-numerator}
\end{align}
For the longitudinal polarization $e=p$ it is
\[
 N_L=(B-A)^2\left[D_s-2+\frac{(A+B-1)^2}{4AB}\right].
\]
The $D_s$ dependence comes from the identity part of the vector
contraction. One added direction orthogonal to $e,p,k$ contributes
$(2e\cdot k+e\cdot p)^2$, which is $4t^2$ in the transverse projection.
It is nonzero, for example, when $p=(1,0,0,0)$,
$e=(0,1,0,0)$ and $k=(0,1,0,0)$.

For any scalar function of $A,B$, rotational reduction around $p$ gives
\[
 t^2\ \rightsquigarrow\
 \frac{1}{D-1}\left[A-\frac{(B-A-1)^2}{4}\right].
\]
At $p^2=1$, write $I_{ab}=\mu_{\rm DR}^{2\epsilon}
\int d^Dk/(2\pi)^D\,A^{-a}B^{-b}$. The standard Gamma-function formula
gives
\[
 I_{21}=I_{12}=-(D-3)I_{11},\qquad
 I_{22}=(D-6)(D-3)I_{11}.
\]
Terms with a nonpositive denominator exponent are scaleless and vanish
in this dimensional prescription. Some intermediate integrals have
infrared poles; the reduction is an identity of meromorphic integrals,
and their sum is formed before expansion at $\epsilon=0$.

Restoring $p^2$, the result is
\begin{align}
 G^{ab}_{\mu\nu}
 &=g^2 C^{ab}p^2 I_{11}(p)
 \left[R_T(D,D_s)P^T_{\mu\nu}-\frac14P^L_{\mu\nu}\right],
 \label{eq:gluon-reduced}\\
 R_T(D,D_s)&=\frac{7D^2-15D+14-4D_s}{8(D-1)},\notag\\
 I_{11}(p)&=\frac{e^{(L+\gamma_E)\epsilon}}{16\pi^2}
 \frac{\Gamma(\epsilon)\Gamma(1-\epsilon)^2}
 {\Gamma(2-2\epsilon)}
 =\frac{1}{16\pi^2}\left[\frac1\epsilon+L+2+O(\epsilon)\right].\notag
\end{align}
In particular,
$R_T(4-2\epsilon,4-2\epsilon)=25/12-61\epsilon/36+O(\epsilon^2)$.
After subtracting the pole,
\begin{equation}\label{eq:gluon-finite}
 \boxed{G^{ab}_{\mu\nu,R}
 =\frac{g^2C^{ab}p^2}{16\pi^2}
 \left[\left(\frac{25}{12}L+\frac{89}{36}\right)P^T_{\mu\nu}
 -\left(\frac14L+\frac12\right)P^L_{\mu\nu}\right].}
\end{equation}

\subsection{Finite trace contribution and the full two-point check}

Keeping $D$-dimensional integration but replacing $D_s$ by $4$ changes
the transverse coefficient by
\[
 R_T(D,D)-R_T(D,4)=-\frac{D-4}{2(D-1)}.
\]
Multiplication by the pole of $I_{11}$ therefore leaves
\begin{equation}\label{eq:gluon-extra-finite}
 \boxed{G^{ab}_{\mu\nu,R}\big|_{D_s=D}
       -G^{ab}_{\mu\nu,R}\big|_{D_s=4}
 =\frac{g^2C^{ab}}{16\pi^2}\frac13
 (p^2\delta_{\mu\nu}-p_\mu p_\nu).}
\end{equation}
The frozen-trace gluon constant would be $77/36$, instead of $89/36$.
Its pole is unchanged. Its longitudinal part is also unchanged, so
transversality alone will not detect this error. This is an instance of
the finite contributions from dimension-dependent numerators and UV
poles discussed in~\cite{rational}.

The previous ghost calculation, with $g^2$ restored, has transverse
coefficient $L/12+2/9$ and longitudinal coefficient $L/4+1/2$.
Adding it to \eqref{eq:gluon-finite} cancels the longitudinal part and gives
\begin{equation}\label{eq:full-two-point}
 (G+B)^{ab}_{\mu\nu,R}
 =\frac{g^2C^{ab}}{16\pi^2}
 (p^2\delta_{\mu\nu}-p_\mu p_\nu)
 \left[\frac{13}{6}L+\frac{97}{36}\right].
\end{equation}
The massless four-gluon tadpole vanishes as a scaleless integral in
conventional dimensional regularization. Equation~\eqref{eq:full-two-point}
is therefore the complete conventional pure-gauge one-loop two-point
answer in these conventions. A retained cutoff can give a nonzero
tadpole before limits; its source-level treatment is a separate check.

The sum agrees with Grozin's exact-$D$ expression in~\cite{grozin},
equation (84), with $n_f=0$ and $\xi=1-a_0=1$ (Landau gauge).
His equations (78)--(81) place the self-energy with a minus sign in
the inverse propagator, consistent with our log-density convention.
His logarithm is $-L$. The finite coefficient is also listed in
Appendix A.1 of~\cite{ruijl}, with the opposite sign in that paper's
self-energy convention. These are independent comparison references.

\subsection{What has been checked against the source}

The fresh Lean module \code{GluonPolarizationAudit.lean} proves five
diagnostic identities. Two use the actual \code{Landau.precision}:
its trace is $(\#\mathrm{axes}-1)q^2$, and division by $q^2\ne0$ gives
the transverse projector trace. For a fixed color and a two-leg metric
test vertex, the next two prove
\[
 \operatorname{yeast}=1,\qquad
 \operatorname{innerTrace}(z)=(4+2m)+2(z-m)=4+2z
\]
in every sector $m$. The final identity proves, for arbitrary inputs,
\[
 \operatorname{innerTrace}(z)-\operatorname{innerTrace}(z')
 =2(z-z')\,\operatorname{yeast}.
\]
The metric test is a test of the trace operation; it is not asserted to
be an admissible loop integrand. All five proofs use only the standard
Lean axioms recorded in the receipt.

The same run inspects four submitted theorems. \code{field_W} identifies
the actual cubic vertex on physical gauge legs; \code{gauge_hit} and
\code{trace_value} give the actual physical-sector sewing and trace.
\code{closeG_integer} identifies integration in fitting integer sectors
under its admissibility hypotheses. Inspecting these statements is not
an independent proof of their analytic content, and the integer-sector
theorem does not characterize a continuation near $z=0$ by itself.

The independent exact-arithmetic script expands the two projectors and
cubic tensors, reduces the resulting Laurent polynomials in $A,B$, and
obtains \eqref{eq:gluon-reduced}. It also checks 48 explicit-index
contractions in dimensions four, six and eight, 24 added-axis differences,
and 48 comparisons between test sectors. These are checks of an
independent transcription, not execution of the Lean graph evaluator.
The scalar master-integral identities are additionally checked against
their Gamma formulas at six noninteger dimensions. The script verifies
the exact-$D$ sum against Grozin and the finite negative control in
\eqref{eq:gluon-extra-finite}.

\emph{Status.} The full source comparison is now proved. The theorem
\code{GluonSolution.acceptance} establishes \eqref{eq:gluon-finite} for
\code{(TwoPoint.bubble G false false).instructions.renorm}, with the coupling
factor stripped from the formal target. It covers all nonzero physical
momenta, tensor components, colors and positive $\bar\mu$, under the submitted
simple-group hypothesis. The proof starts from the actual cubic sewing and
\code{yeast}, evaluates the Gaussian and radial integrals, differentiates
the continued top Taylor term and fixes the finite integration constant
from the source's ultraviolet subtraction. The unit-probe retained-cutoff
identity is also proved. The separate note
\code{gluon-bubble-finite-identity.pdf} gives the calculation and proof map.

A fresh review build and frozen-target Comparator check accept the theorem
with only the standard Lean axioms and no permitted definition holes.
This proves the stated graph comparison within the submitted definitions;
it does not complete the human review or establish the conventional
all-orders reference identification. The finite $1/3$ contribution remains
a useful diagnostic: checking the pole or transversality alone would miss it.

\section{Review of the seven convention checks}
\label{sec:completion}

This section records an automated source and mathematical review of every
row in Section~\ref{sec:dictionary}. It adds no human endorsement. The
received source was checked against its manifest, and the existing review
receipts were checked for integrity. The new Lean check proves a simplification
of the actual cubic density, unfolds the quartic density, and inspects nine
submitted bridge theorems. Independent exact algebra tests check the gauge
vertices and the leading Wilson plaquette coefficient. They do not execute
the Lean graph evaluator or reprove the general matching theorem.

\subsection{Color and coupling}

The color convention is explicit. In \code{YM/Gauge/Colors.lean},
\code{colorIso} rescales a Frobenius-orthonormal frame by
$\sqrt{d_{\rm rep}}$, and \code{trace_color} proves
\[
 -\Re\tr(E_aE_b)/d_{\rm rep}=\delta_{ab}.
\]
The bracket is the matrix commutator transported through that frame.
\code{trace_adjoint}, under the simplicity hypothesis, proves
$-\Tr(\ad_x\ad_y)=C_A\langle x,y\rangle$. These are identities about
the actual color space, rather than a separate choice of structure constants.

For SU($N$) in its fundamental representation, comparison with the common
anti-Hermitian generators $t_a$ normalized by
$-\tr(t_at_b)=\delta_{ab}/2$ gives $E_a=\sqrt{2N}\,t_a$, up to an
orthogonal change of basis. Thus $f_E=\sqrt{2N}\,f_t$,
$C_A^{E}=2N^2$, and the canonical interaction couplings obey
$g_{\rm std}=\sqrt{2N}\,g$. In particular SU(2) has $C_A^{E}=8$,
consistent with the previously checked Pauli basis. A different chosen
orthonormal frame does not change this invariant. This concerns the formal
action coupling; $g_X$ additionally uses the prescribed finite conversion.

\emph{Assessment.} The normalization issue is resolved at the level needed
to state and compare these rules. Matching a particular coordinate frame
can still be formalized, but it is ordinary algebra. The interpretation of
the measured coupling is discussed below; it should not be confused with a
missing color normalization.

\subsection{Gauge covariance}

The existing Lean identities compare the actual unfiltered line and twice
the colored half-line with the explicit transverse propagator. The factor
of two is accounted for by the second-order Wick operator. Independently,
contracting the displayed tensor with $p_\mu$ gives zero. With
$u=e^{2s}p^2$, the filter satisfies
\[
 \partial_s[1-(1+u)e^{-u}]=2u^2e^{-u}.
\]
This also checks the sign and normalization of the shell derivative; the
derivative of the half-line carries $u^2e^{-u}C_{\mu\nu}(p)$.

\emph{Assessment.} No substantial interpretive ambiguity remains in the
free momentum-space rule at $p\ne0$. The value assigned by a totalized
inverse at $p=0$ is not a distributional prescription. Position-space
extension, zero modes at finite volume, and interchanges of loop integration
with cutoff limits are separate analytic tasks. They can be specified and
checked mathematically; their difficulty alone is not a reason to classify
them as requiring human interpretation.

\subsection{Ghost orientations and the source restriction}

The existing five Lean identities fix both ordered propagators, the cubic
vertex of the source action, and its scalar-antighost form. They use
$K_\mu(q)=-iq_\mu\bar c(q)$, the derivative convention $+ip_\mu$, and
the stated external order. Together with the actual two-point ghost-loop
comparison, these are substantial checks of the signs and normalization.

The restriction of $K$ is nevertheless a restriction. The unrestricted
vector antifield has more components than a scalar antighost, and
recovering an antighost at zero momentum would require additional data.
Neither the local identities nor the two ghost-loop examples prove an
equivalence of all renormalized BV source insertions with a conventional
gauge-fixed operator system. No global nonperturbative gauge-fixing claim
is needed for the local perturbative dictionary being tested here.

\emph{Assessment.} The elementary ghost rules are checked. Human review
should concentrate on the physical meaning of the source restriction when
composite insertions, equations of motion, and BRST-exact terms occur. A
general theorem transporting those operations through gauge fixing remains
a mathematical task, rather than something a reviewer should accept on
the strength of notation.

\subsection{Cubic and quartic gauge vertices}

\code{Class_at} identifies the interaction seeds with the full classical
log density minus its free part. The new Lean identity simplifies
\code{top3} from its symmetrized form to $-F_0Q/2$ using the even parity
of the gauge fields. The quartic seed is exactly $-Q^2/4$. The source's
shuffle product allocates external derivative slots to factors, so ordinary
differentiation of these polynomials supplies the vertex multiplicities.

Writing $f^{abc}=f^{ab}{}_c$ in an orthonormal color frame, differentiation
gives the following log-density tensors, with all momenta incoming:
\begin{align}
 V^{abc}_{\mu\nu\rho}(p,q,r)
 &=ig f^{abc}\bigl[
 \delta_{\mu\nu}(p-q)_\rho+
 \delta_{\nu\rho}(q-r)_\mu+
 \delta_{\rho\mu}(r-p)_\nu\bigr],\label{eq:checked-cubic}\\
 V^{abcd}_{\mu\nu\rho\sigma}
 &=-g^2\sum_e\Bigl[
 f^{abe}f^{cde}(\delta_{\mu\rho}\delta_{\nu\sigma}
                  -\delta_{\mu\sigma}\delta_{\nu\rho})\nonumber\\
 &\qquad+f^{ace}f^{bde}(\delta_{\mu\nu}\delta_{\rho\sigma}
                  -\delta_{\mu\sigma}\delta_{\nu\rho})\nonumber\\
 &\qquad+f^{ade}f^{bce}(\delta_{\mu\nu}\delta_{\rho\sigma}
                  -\delta_{\mu\rho}\delta_{\nu\sigma})\Bigr].\label{eq:checked-quartic}
\end{align}
For the cubic identity, allocating three distinct derivative slots to
$F_0Q$ gives six assignments. For the quartic identity, allocating four
slots to $Q^2$ gives twenty-four. Antisymmetry and cyclicity of $f$ collect
these assignments into the displayed tensors. No further $3!$ or $4!$
divisor is present.

An independent exact-arithmetic implementation compares these unsimplified
assignments with the displayed expressions. It checks all Lorentz
components for two alternating color tensors of dimensions three and four:
69,888 cubic momentum coefficients and 86,272 quartic tensor entries.
The cubic check treats all external momenta as independent, before imposing
momentum conservation. Nonzero entries detect a reversed overall sign or
an extra vertex factorial. The color-dimension-four case is an algebraic
test tensor, not a claim about a new simple gauge group. Four noncommuting,
affine SU(2) plaquette tests provide an additional normalization check below.

\emph{Assessment.} The classical gauge vertices now have a written algebraic
comparison and exhaustive finite component tests in the stated cases.
A single generic Lean theorem equating every submitted three- and four-leg
row to these tensors has not been proved here. Completing that theorem is
mechanical work. It is not a reason to leave the underlying convention
undecided or to require expert interpretation of the elementary vertex.

\subsection{Subtraction and dimensional continuation}

The source does more than discard divergent terms from individual
integrands. \code{integrateRen} includes its children's counterterms,
integrates the resulting scale sector, and inserts the negative principal
part of the resulting Laurent germ. In the explicit history representation,
\code{GraphSum.forest_germ} proves equality with the recursively enumerated
quotient-graph sum, retaining the dimensional parameter in the contracted
complements. The theorem is for the source's scale-ordered forest
prescription. It is not itself a comparison with an independently specified
incidence-graph forest prescription.

The relation $D=4+2z=4-2\epsilon$ and the shift
$y=\log\bar\mu+(\gamma_E-\log4\pi)/2$ give the stated
$\MS$ loop-measure convention. The explicit continuation contains both
a split radial integral and an extra-polarization trace:
\[
 \mathrm{innerTrace}=\hbox{trace over the test axes}
                   +2(z-m)\,\mathrm{yeast}.
\]
Ordinary integral identities at integer dimension parameters check its
normalization. They cannot alone characterize the continuation near zero.
The source itself specifies this noninteger-dimensional rule, including
the continuation of open tensor slots.

The two ghost examples test finite subtraction constants, but their closures have no extra vector-polarization trace. Section~\ref{sec:gluon-bubble} now identifies the submitted gluon bubble through its finite endpoint in Lean and Comparator. It includes a finite $1/3$ contribution from the continued vector trace that both a pole-only check and transversality would miss. A two-loop example with overlapping divergent subgraphs would also test an aspect absent from these primitive examples.

\emph{Assessment.} The source is a definite local pole-subtraction
construction, with real evidence beyond a definition bearing the name MS.
The remaining all-order comparison with conventional dimensional
renormalization is not finished. Human review is particularly valuable in
fixing what ``the conventional $\MS$ reference'' means for these open
tensors and composite operators. Evanescent changes of an operator lift
can combine with pole counterterms to change finite coefficients. Once
the comparison prescription is stated, forest equivalence and the
additional loop tests are mathematical tasks that can be pursued without
waiting for a human to perform the calculations.

\subsection{Local fields, composite insertions, and OPE}

The public labels are gauge-invariant free on-shell curvature-jet
polynomials. The graph correlation expands a label in a finite basis and
uses the actual insertion polynomial chosen by \code{OpenAxis.Plot}.
The source then constructs the marked kernels and their counterterms.
This is more specific than assigning an arbitrary perturbative series to
each label.

A stronger internal identification is available than the previous table
indicated. In \code{ScaleForest/Optics.lean}, \code{common_label},
\code{common_column} and \code{common_native} identify the OpenAxis and
JRad insertion constructions, including their auxiliary-dimensional
components, before applying measured field frames. Thus agreement is not
merely equality of four-dimensional labels after forgetting evanescent
data. The corresponding uniqueness theorems identify the marked Laurent
germs on these inputs.

For the public OPE, \code{HasChannelBasis} requires the channels through
the output cutoff to form the whole finite field space. \code{measuredOPE}
applies the inverse measured frame to the summed physical output.
The estimate tests the full coefficient norm and the dual unit ball;
it is not just an equality of vacuum expectations. The reference OPE is
characterized by arbitrary exterior insertions.

\emph{Assessment.} The internal label and insertion bridges are stronger
than ``definitions located.'' The unresolved external comparison concerns
the meaning of the renormalized operators. In particular, a free label
such as $F_0^2$ does not by its name identify a conventional interacting
normal product $[F^2]_{\MS}$: the covariant completion, mixing, equation-of-motion
terms and continuation must be compared. Human review should decide whether
the claimed physical operator system is the intended one, and whether the
separated-point scope is made explicit. Contact identities and an arbitrary
off-shell operator algebra do not follow from this OPE statement.

\subsection{Wilson normalization and calibration}

The public lattice law has precisely the normalized Wilson density
$\exp(-\beta\sum_p[1-\Re\tr(U_p)/d_{\rm rep}])$ against product Haar
measure, with each positively oriented plaquette counted once.
\code{energy_eq}, \code{law_eq} and \code{moment_eq} in
\code{WilsonBridge.lean}
identify the public formulas with the internal energy, measure and finite
moments. \code{AdmissibleLattice} also retains the native insertion
coordinates. These checks rule out a discrepancy at this public interface;
they do not by themselves check the full renormalization construction.

Independently, for a smooth anti-Hermitian connection $\mathcal A$,
the plaquette has logarithm $a^2\mathcal F_{\mu\nu}+O(a^3)$. Its trace gives
\[
 1-\Re\tr U_p/d_{\rm rep}
  =\frac{a^4}{2}|\mathcal F_{\mu\nu}|^2+O(a^5).
\]
Summing $\mu<\nu$ produces $\tfrac14\sum_{\mu,\nu}\mathcal F_{\mu\nu}^2$.
Writing $\mathcal A=gA$ and $\beta=g^{-2}$ therefore gives the canonical
action used above at leading order. The exact finite tests multiply the
four link exponentials through degree four for affine noncommuting SU(2)
fields. Their $a^4$ coefficients are $2$, $71$, $7$, and $129/2$, each equal
to $|\partial_\mu\mathcal A_\nu-\partial_\nu\mathcal A_\mu+
[\mathcal A_\mu,\mathcal A_\nu]|^2/2$. This checks the trace, orientation,
derivative and commutator factors in nontrivial cases; it is not an estimate
uniform over fluctuating lattice fields.

The public calibration is also restrictive. Its coupling is an exact signed
triangle measurement divided by a fixed $c_T$, rather than an adjustable
parameter chosen separately for each coefficient. The field frame solves
fixed position normalization conditions. The reference is fixed before
the hierarchy, and a hierarchy's frame and coupling are fixed before the
perturbative order. \code{Reference.scheme_unique} proves uniqueness for
the prescribed group menu. Higher-order tuning preserves earlier orders
in \code{High/Drive.lean}; the actual tuning data are supplied separately.

\emph{Assessment.} The elementary Wilson normalization and public measure
identification are checked. A full independent audit of the lattice flow,
implemented local tuning and remainder bounds remains technical work.
The principal interpretive point is the scope of the conclusion: it is
matching in these calibrated fields and this measured coupling, after the
fixed finite conversion from MS. It does not assert equality for arbitrary
preassigned bare field coordinates. Expert review should assess that
scientific claim and the allowed operator mixing, rather than simply
rechecking which variable is named ``coupling.''

\subsection{Where expert judgment should be concentrated}

Three questions deserve priority. First, does the explicit continuation
and subtraction prescription coincide with the intended conventional
reference, or is a further finite conversion required? Second, do the
continued insertion representatives, their BRST interpretation and their
equation-of-motion quotient describe the physical composite operators
claimed in the assessment? Third, is the calibrated, separated-point
matching theorem the intended scientific claim, with its field changes,
order of limits and restrictions on contact terms stated accurately?

These are questions about the interpretation and reach of a precise
formal result. They are not evidence of a counterexample. Nor are they
intrinsically beyond formal analysis: a sufficiently explicit external
specification would turn much of this work into further comparison
theorems. In contrast, a generic tensor identity, a graph enumeration,
a two-loop benchmark, a source-to-source equality or a limit estimate
should remain listed as a technical task, even when it is substantial.
This review found no mismatch in the checks carried out. It has not
completed an independent all-order equivalence proof.

\section{Review priorities: code coverage and spot checks}
\label{sec:review-plan}

The priority is to test the subtraction prescription and the meaning of
the composite insertions. A spot check should exercise a convention that
could change an answer, compare with a separately stated reference, and
report its limited scope. Passing selected examples cannot establish
equivalence for all operators and all orders. The present plan leaves that
general confirmation to future work.

\subsection{Existing code for the subtleties}

The submission contains explicit definitions and theorem bodies for the
following mechanisms. Locating them is evidence that the issues have been
addressed within the construction; it is not an independent verification
of their conventional interpretation. Helper theorems must be read with
their hypotheses, and their use for a concrete benchmark must be traced.

\begin{itemize}
\item \textbf{Subdivergences and dimensional continuation.}
\code{YM/MS/Graphical/Block.lean} defines \code{selectCT} and
\code{integrateRen}; \code{ScaleForest/Quotients.lean} defines the recursive
contracted-history sum and proves \code{forest_germ}. The complement retains
the regulator. \code{ScaleForest/Glue.lean} defines \code{innerTrace} and
\code{closeG}, including the extra-dimensional trace and split radial
integral. The conventional identification of overlapping incidence-graph
subtractions is still an external comparison obligation.
\item \textbf{Nonlinear insertions and their continued representatives.}
\code{YM/FVA2/Gen/Recipe.lean} relates covariant curvature jets to the radial
jet construction, including interactions. \code{Position/Table.lean}
constructs \code{OpenAxis.Plot}; \code{YM/MS/ScaleForest/Optics.lean}
proves the common-label and common-insertion bridges in the continued
setting. \code{ChooseRad.lean} gives uniqueness within its specified class
of lifts. It does not identify arbitrary alternative evanescent prescriptions.
\item \textbf{Mixing, exact directions and the equation quotient.}
\code{YM/MS/Local/Physical.lean} defines \code{Null} using the image of
\code{natQ}; \code{High/Adjust.lean} proves a physical-plus-exact
decomposition, and \code{High/Standard.lean} uses it in marked tuning.
\code{YM/Maxwell/Equation.lean} implements the free on-shell substitution
\code{reduceEq} and characterizes its kernel. These are concrete mechanisms;
identification with conventional interacting BRST and equation-of-motion
relations, especially their contact terms, needs a separate comparison.
\item \textbf{Reduction to physical operators and scheme conversion.}
\code{YM/FVA2/Extras/Result.lean} proves \code{certified_reduction} under
full-dimensional closure and admissibility hypotheses;
\code{Extras/Proj.lean} proves \code{rewrite_extra} in separated correlations.
The source explicitly excludes comparison with other dimensional tensor
definitions of evanescents from this result. \code{MSConversion.lean}
specifies the field and coupling series; the public OPE uses the inverse
frame on its output. None of these statements by itself supplies an
independent conventional mixing matrix for a chosen operator.
\end{itemize}

\subsection{Prioritized independent checks}

The following are proposed checks, not newly completed results. Begin with
the first two; use the others to broaden coverage as time permits. The
existing gluon-bubble finite equality already tests an evanescent trace
effect, so another primitive bubble would add less coverage.

\begin{enumerate}
\item \textbf{One scalar operator through the full dictionary.}
Choose a definite normalization of $\mathcal O=\tr F_{\mu\nu}F_{\mu\nu}$.
Trace its public free label to its interacting insertion. Check the
two-, three- and four-gauge-field tree vertices, using noncommuting colors
so that the commutator terms are visible. Compare one one-loop marked
amplitude, including its pole, finite part and allowed mixing, in explicitly
matched conventions. Use nonexceptional external momenta and state the
connected/amputated or 1PI convention. If the off-shell marked kernel is
not directly exposed by the public reference, state and verify that bridge.
This tests one operator's completion and normalization, not the whole algebra.
\item \textbf{A genuine subdivergent subtraction.}
Select an actual two-loop Yang--Mills history with a proper divergent
one-loop subgraph. Compare the subdivergence counterterm, its insertion
into the regulator-dependent complement, the remaining overall pole, and
one finite projection with an independent forest calculation. Record any
IR regulator needed and its removal. A nested example tests recursion;
it does not test overlaps. A separate overlapping-divergence example is
the next extension, not an automatic consequence of the nested test.
\item \textbf{A derivative and a redundant direction.}
Check a derivative descendant of the scalar insertion, including the
Fourier momentum factor and compatibility with subtraction. Trace one
explicit equation-of-motion or exact insertion through its actual closure
and reduction hypotheses, then compare its effect in separated physical
correlations. Do not assert that every off-shell amplitude vanishes:
contact terms and the distinction between free and interacting equations
must be retained in the comparison.
\item \textbf{An evanescent change of representative.}
For concrete admissible insertions with the same quantum numbers, compare $\mathcal O_D$ with
$\mathcal O_D+\epsilon\mathcal P_D$ using the same convention
$z=-\epsilon$. Track the finite shift induced by a pole of
$\mathcal P_D$ through the actual subtraction and physical reduction.
This is a controlled sensitivity test, not a claim of scheme independence.
More general dimension-specific tensor identities remain future work.
\item \textbf{One finite conversion and OPE output.}
For the chosen scalar sector, calculate the first nontrivial field/coupling
conversion and verify one OPE coefficient with the inverse frame on the
output. Keep all channels allowed by the selected finite cutoff. This
tests that an agreement in MS is transported to the stated measured scheme.
\end{enumerate}

Each completed check should record the pinned source definitions, external
reference conventions, independent expected value, actual source expression,
and whether the comparison is analytic, numerical, or proved in Lean and
Comparator. Negative controls should detect the specific error being tested,
such as dropping a child counterterm, omitting the commutator completion, or
setting $D=4$ before pole subtraction. A check whose reference convention or
source bridge is still unresolved remains pending.

The appropriate assessment is therefore: selected finite amplitudes agree
with independent conventional calculations, and explicit source machinery
addresses the broader conventions; exhaustive identification of subtraction
and renormalized composite operators has not been established by this
review. The all-orders flow and composite-operator literature~\cite{fhh,fh} supplies
comparison frameworks, not a blanket certificate that its conventions
agree with this implementation.

\appendix
\begingroup\small
\section{Source references and verification}
\label{sec:source}

The source is commit \code{b5e288edeac43ac088a2a291ac3841bdda9f20b8} of
\href{https://github.com/yang-mills-group/OpenAI-YM-20261104}{YM-lean/OpenAI-Nov4}.
Paths below are relative to \code{submission/lean/}.
\code{assessment-source-map.json} records file hashes and declaration
locations. The longer \emph{Perturbative matching in the November 4 OpenAI
Yang--Mills construction} notes contain the detailed audit.

\begingroup\small
\begin{longtable}{@{}>{\raggedright\arraybackslash}p{0.25\textwidth}>{\raggedright\arraybackslash}p{0.70\textwidth}@{}}
\toprule Topic & Files and declarations\\\midrule
\endfirsthead
\toprule Topic & Files and declarations\\\midrule
\endhead
Public target & \code{YangMills/PerturbativeMatching.lean}: \code{MatchedRun}, \code{MatchedRefinement}. Its subdirectory contains \code{Correlations}, \code{OperatorProducts}, \code{Measurement}, \code{Reference} and \code{Asymptotics}.\\[4pt]
Common flow algebra & \code{YM/MS/CoverFlow.lean}: \code{using_tree}, \code{recu}, \code{Track}, \code{Flow}, \code{recu_cost}. \code{CostVertex.lean}: \code{cycle} and comparison norms; \code{CompareSymbol.lean}: weighted difference.\\[4pt]
Lattice versus continuum RHS & \code{YM/MS/LandauMatch.lean}: \code{zL}, \code{zD}, \code{land_value}. \code{YM/MS/High/Sphere.lean}: \code{Spheres}, \code{Trust}, \code{first_rhs}, \code{second_rhs}.\\[4pt]
Local obstruction and tuning & \code{YM/MS/High/Iterate.lean}: \code{advanced}, \code{closes}, \code{terminal}. \code{High/Standard.lean}: \code{action}, \code{mark}. \code{High/Drive.lean}: \code{Fixed}, preservation of earlier choices.\\[4pt]
Actual comparison inputs & \code{YM/MS/Wind/Traces.lean}, \code{Wind/Clone.lean}, \code{Wind/Portfolio.lean}: lattice flow and \code{portH}. \code{YM/MS/Image/Verdant.lean}: \code{Outlook.dock}, \code{fieldD}.\\[4pt]
MS solution and installation & \code{YangMills/WilsonProof/NativeSupport/YM/MS/Mach/Land.lean}: \code{MS_Data}, \code{leader_M}, \code{High_stops}, \code{uniq}. \code{YM/WilsonMS/Calibration/StrongWind.lean}: \code{Installed}, \code{can_install}.\\[4pt]
Graph/flow identity & Under \code{YangMills/WilsonProof/NativeSupport/YM/MS/Graphical/}: \code{Leading.lean} has \code{ambient_bob}, \code{mark_bob}, \code{income_bulk}, \code{income_mark}; \code{Diagram.lean} explains sectors; \code{Block.lean} proves subtraction properties. The definitions of \code{Diagram} and \code{integrateRen} are in the corresponding \code{YM/MS/Graphical/} files.\\[4pt]
Ghost-box source rules & Under \code{YM/MS/ScaleForest/}: \code{Ordered.lean}, \code{Glue.lean}, \code{AllHistories.lean}, \code{Reaching.lean}, \code{Entry.lean}, \code{Halves.lean}, \code{HeatScale.lean}, \code{TwoVertices.lean} and \code{PageOne.lean}. Koszul allocation is in \code{YM/MS/KoszulAllocate.lean}; the radial measure is in \code{YM/MS/Train/MeasureD.lean} and \code{YM/MS/Arc/Pumps.lean}.\\[4pt]
Gluon trace and sewing & \code{YM/MS/Landau.lean}: \code{precision}. \code{YM/MS/Arts/Tall.lean}: \code{yeast}. \code{YM/MS/ScaleForest/Glue.lean}: \code{innerTrace}, \code{closeG_integer}. \code{YM/WilsonMS/Beta/Flash.lean}: \code{field_W}. \code{ScaleForest/TracedCubic.lean}, \code{LoadB.lean}, \code{Bubbles.lean}: \code{gauge_hit}, \code{trace_value}, \code{bubble}.\\[4pt]
Coefficient limits & \code{YM/FVA2/XReg/Continuum.lean}; \code{YM/WilsonMS/Scales/PointMove.lean}. Under \code{YM/Integration/FVA1Physical/Exchange/}: \code{ContinuumLimitThermometer.lean} and \code{CalibratedExpansionCommonCalibratedGates.lean}.\\[4pt]
Graphs to physical bounds & \code{YM/MS/Graphical/CommonGates.lean}: \code{calibrated_gates}. \code{YM/SignedWilsonExpansion/Result/Orders.lean}: \code{pricedLines}. \code{YM/Integration/WilsonGraphMeasuredCompactCorrelation.lean}: \code{publicMixedPolynomial}.\\
\bottomrule
\end{longtable}
\endgroup

In the verification repository, the existing checks are
\code{checks/matching/AllOrdersStatement.lean},
\code{GhostConventionCheck.lean}, and the
\code{GhostLoopSolution.lean} proof with its supporting modules.
\code{docs/matching/ghost-finite-check.md} states the benchmark's exact
scope. Its receipts are in \code{evidence/matching/ghost-finite-build.json}
and \code{evidence/matching/ghost-finite-comparator/}. These checks ran in
macOS development mode.

The independent conventional box calculation is reproduced by
\code{scripts/check_ghost_box_reference.py}; its output is
\code{evidence/matching/ghost-box-reference.json}.
The new source check is \code{checks/matching/GhostBoxSource.lean}, with
build and axiom receipts in \code{evidence/matching/ghost-box-source.json}
and its accompanying log. The build reuses the submitted libraries after
checking all 37,673 source files against the pinned manifest; it freshly
compiles the review module. Reproduce it with
\code{scripts/build_ghost_box_source.py --source <built-lean-root> --work <fresh-directory>}.
The independent transcriptions are
\code{scripts/check_ghost_box_sectors.py} and
\code{scripts/check_ghost_box_cutoff.py}, with corresponding JSON records
under \code{evidence/matching/}. Neither executes the submitted Lean
integrals. The box's final finite-value proposition remains unproved in
Lean. The Comparator run recorded in Draft 0.6 concerns the gluon bubble;
Draft 0.7 adds no new formal run.
The section-2.3 review adds \code{checks/matching/Section23Audit.lean}
and \code{scripts/check_section23_algebra.py}. Its build receipt and exact
algebra results are \code{evidence/matching/section23-audit.json} and
\code{section23-algebra.json}; \code{docs/matching/section23-source-map.json}
records the additional source anchors. The Lean run checks two density
identities and reports nine submitted bridge theorems, with standard axioms
only. It does not prove a conventional all-order forest equivalence.
\code{checks/matching/GluonPolarizationAudit.lean} adds five trace identities and inspects four submitted bridges. Reproduce its fresh build with \code{scripts/build_gluon_polarization_audit.py}; the receipt is \code{evidence/matching/gluon-polarization-audit.json}. The independent reference calculation is \code{scripts/check_gluon_bubble_reference.py}, with results in \code{evidence/matching/gluon-bubble-reference.json}. The additional source hashes are recorded in \code{docs/matching/gluon-bubble-source-map.json}. The full gluon proof is in the \code{GluonLoop*.lean} modules, culminating in
\code{GluonLoopSolution.lean}. Reproduce the fresh 33-module build with
\code{scripts/build_gluon_finite.py}, then the frozen-target check with
\code{scripts/run_gluon_finite_comparator.py}. The corresponding evidence is
\code{evidence/matching/gluon-finite-build.json} and
\code{evidence/matching/gluon-finite-comparator/}. The received libraries are
reused; all 37,673 source files are verified before and after Comparator.
The macOS run uses upstream's development wrapper and the Lean default
kernel, without sandbox isolation or an optional second kernel.
\code{docs/matching/assessment-document-checks.json} records compilation,
source hashes and visual inspection of this draft.
\code{docs/matching/convention-review-source-map.json} records the source
anchors inspected for Section~\ref{sec:review-plan}, distinguishing located
mechanisms from completed independent comparisons.

\begin{thebibliography}{9}
\bibitem{srednicki}
M. Srednicki, \emph{Quantum Field Theory}, Cambridge University Press, 2007,
Sections 71--72 (gauge-fixed path integral and nonabelian Feynman rules).
\href{https://web.physics.ucsb.edu/~mark/qft.html}{Author's book page};
\href{https://web.physics.ucsb.edu/~mark/ms-qft-DRAFT.pdf}{author's draft},
equations (71.19)--(71.23) and (72.1)--(72.9).

\bibitem{fp}
L. D. Faddeev and V. N. Popov,
\emph{Feynman diagrams for the Yang--Mills field},
Physics Letters B \textbf{25} (1967), 29--30.
\href{https://doi.org/10.1016/0370-2693(67)90067-6}{doi:10.1016/0370-2693(67)90067-6}.

\bibitem{grozin}
A. Grozin, \emph{Quantum Chromodynamics} (2012), Section 4.2,
equations (78)--(86), \href{https://arxiv.org/abs/1205.1815}{arXiv:1205.1815}.

\bibitem{ruijl}
B. Ruijl, T. Ueda, J. A. M. Vermaseren and A. Vogt,
\emph{Four-loop QCD propagators and vertices with one vanishing external momentum},
JHEP \textbf{06} (2017), 040, Appendix A.1,
\href{https://arxiv.org/abs/1703.08532}{arXiv:1703.08532}.

\bibitem{rational}
J.-N. Lang, S. Pozzorini, H. Zhang and M. F. Zoller,
\emph{Two-Loop Rational Terms in Yang--Mills Theories},
JHEP \textbf{10} (2020), 016,
\href{https://arxiv.org/abs/2007.03713}{arXiv:2007.03713}.

\bibitem{egk}
A. N. Efremov, R. Guida and C. Kopper,
\emph{Renormalization of SU(2) Yang--Mills theory with Flow Equations},
\href{https://arxiv.org/abs/1704.06799}{arXiv:1704.06799}.
See especially the action and flow-equation setup in Sections 2.2--2.3.

\bibitem{fhh}
M. B. Fr\"ob, J. Holland and S. Hollands,
\emph{All-order bounds for correlation functions of gauge-invariant
operators in Yang--Mills theory}, J. Math. Phys. \textbf{57}, 122301 (2016),
\href{https://arxiv.org/abs/1511.09425}{arXiv:1511.09425}.
BRST, composite insertions and flow estimates.

\bibitem{fh}
M. B. Fr\"ob and J. Holland,
\emph{All-order existence of and recursion relations for the operator
product expansion in Yang--Mills theory},
\href{https://arxiv.org/abs/1603.08012}{arXiv:1603.08012}.
An all-order OPE comparison relevant to the scope of the reference.
\end{thebibliography}
\endgroup
\end{document}
